\section{Protecting an arbitrary stack tail with two truth entries}
\label{sec:two-truth-tail}

The new blank initialization reaches the logical source root with
\[
 (N,T,W,A,D,H,X,Y,L)=(1,0,5,0,0,0,0,0,0).
\]
The stack code $5$ is not a stack consisting entirely of theorems:
with the inherited Cantor pairing convention,
\begin{equation}
 P(a,b)=\frac{(a+b)(a+b+1)}2+b,\qquad
 5=P(0,P(0,1)).
 \label{eq:tail-five}
\end{equation}
It contains two truth entries above a tail whose first entry is falsity.
It would therefore be incorrect to invoke the usual invariant that
\emph{every} decoded stack entry is a theorem. We prove instead that the
actual scheduler never reaches that tail.

The proof concerns the current source: command order
$\mathrm{NIP},\mathrm{DROP},\mathrm{MP}$, with fields read only by the
commands that need them. The earlier protected-tail theorem for a common
field reader and order $\mathrm{MP},\mathrm{DROP},\mathrm{NIP}$ does not
supply this conclusion. Both the command order and the field arities
matter below.

\subsection{The exact schedule and its conservative stack balance}

Let $f(n)$ be the second component of the inverse of $P$.
Write a freshly selected candidate as $n=P(t,y)$.
Its first command and remaining proof-list pointer $r$ are
\begin{equation}
\begin{array}{c|c|r|c}
 t&\text{command}&\text{conservative stack change}&r\\ \hline
 0&\mathrm{NIP}&-1&y\\
 1&\mathrm{DROP}&+1&y\\
 2&\mathrm{MP}&0&f(y)\\
 t\geq3&\mathrm{AX}&+10&f(f(y)).
\end{array}
\label{eq:tail-command-arities}
\end{equation}
For MP the first unpairing of $y$ supplies its additive field.
For AX the two unpairings supply its two fields; its remaining
parameter is $c=t-3$. NIP and DROP have no additional field.
If $r>0$, continuation fetch decrements it before decoding, so the next
packed command is $r-1$. If $r=0$, the candidate is finished.

Define a finite command word $\mathcal W(n)$ recursively by taking the
first command in \eqref{eq:tail-command-arities}, followed by
$\mathcal W(r-1)$ when $r>0$. The pairing inequalities give
$r\leq y\leq n$, so this recursion is well founded. In particular,
\begin{equation}
 |\mathcal W(n)|\leq n+1.
 \label{eq:tail-word-length}
\end{equation}
Indeed, a nonempty recursive suffix has length at most $r$ by induction.
A nonhalting source run from $N=1,L=0$ executes
$\mathcal W(1),\mathcal W(2),\ldots$ in that order. Registers containing
formulas affect MP success and possible halting, but do not change this
schedule before a halt.

The conservative change is exact for DROP, AX and a nonhalting NIP.
A successful MP pushes a truth entry, while an unsuccessful MP leaves
the stack alone; assigning MP the value zero is therefore a lower
bound. AX pushes the old top, one generalized cache formula, and eight
logical rows, irrespective of whether its theory-bundle gate is open.
A halting NIP is treated as having change $-1$ only in the abstract word:
the actual run has already stopped before its pop.

Let $b(n)$ be the total conservative balance of $\mathcal W(n)$ and let
$m(n)$ be its minimum prefix balance, including the empty prefix.
If the first-command change is $\delta$ and its suffix pointer is $r$,
then
\begin{align}
 b(n)&=\delta+\begin{cases}b(r-1)&r>0,\\0&r=0,\end{cases}\notag\\
 m(n)&=\min\left(0,\delta+
               \begin{cases}m(r-1)&r>0,\\0&r=0.\end{cases}\right).
 \label{eq:tail-recursions}
\end{align}
Set $B(-1)=0$ and $B(n)=\sum_{i=0}^n b(i)$.
Since every command has change at least $-1$,
\begin{equation}
 m(r-1)\geq-r\qquad(r>0).
 \label{eq:tail-suffix-loss}
\end{equation}
We use the same bound with value zero for an empty suffix.

\begin{theorem}[Exact conservative debt]
\label{thm:tail-exact-debt}
Every command prefix of
$\mathcal W(0)\mathcal W(1)\mathcal W(2)\cdots$ has conservative
balance at least $-3$. Every command prefix after omitting
$\mathcal W(0)$ has balance at least $-2$, and this latter lower bound
is attained. Furthermore, $B(n)\geq7$ for every $n\geq6$.
\end{theorem}

The assertion is about this conservative accounting. It does not claim
that two initial truth entries are necessary for every actual execution:
successful MP commands can supply extra entries.

\subsection{An exact finite base}

Write $\Delta_s=s(s+1)/2$. Cantor diagonal $s$ consists, in order, of
\[
 P(s-y,y)=\Delta_s+y,\qquad 0\leq y\leq s.
\]
The base covers the complete diagonals $0,\ldots,35$, namely the 666
candidates $0,\ldots,665$ because $\Delta_{36}=666$.

\begin{lemma}[Finite base]
\label{lem:tail-finite-base}
For $0\leq n\leq665$,
\[
 B(n-1)+m(n)\geq-3.
\]
For $6\leq n\leq665$, $B(n)\geq7$, and $B(665)=5762$.
The first six values of $B(n)$ are
$-1,0,-2,-2,-2,-2$.
\end{lemma}

\begin{proof}
These statements follow by exact integer evaluation of
\eqref{eq:tail-recursions}. The following complete algorithm specifies
the calculation. All array references on its right sides are to
smaller indices, as proved above.
\begin{verbatim}
b = []; m = []; length = []; B = []
total = 0
for n in range(666):
    t, y = unpair(n)
    if t == 0:   delta, r = -1, y
    elif t == 1: delta, r =  1, y
    elif t == 2: delta, r =  0, second(y)
    else:        delta, r = 10, second(second(y))
    assert r == 0 or r - 1 < n
    b.append(delta + (b[r-1] if r else 0))
    m.append(min(0, delta + (m[r-1] if r else 0)))
    length.append(1 + (length[r-1] if r else 0))
    assert length[n] <= n + 1
    assert total + m[n] >= -3
    total += b[n]
    B.append(total)
assert min(B[6:]) == 7
assert B[665] == 5762
assert B[:6] == [-1, 0, -2, -2, -2, -2]
\end{verbatim}
Here the inverse is evaluated without floating point:
if $d=\lfloor(\sqrt{8n+1}-1)/2\rfloor$ and $e=n-\Delta_d$,
then $\operatorname{unpair}(n)=(d-e,e)$.
The supplied independent checker uses integer square root.
Its loop invariant is exactly the definition of $b,m,B$ and
\eqref{eq:tail-word-length}; thus this is a finite proof of the induction
base, not a bounded replacement for the infinite assertion.
\end{proof}

For example, the first ten total word balances are
\[
 -1,\ 1,\ -2,\ 0,\ 0,\ 0,\ 10,\ 0,\ 2,\ -3.
\]
Candidate 5 begins with NIP at global balance $B(4)=-2$,
so its first command attains $-3$.

\subsection{Two prefix-sum transforms under unpairing}

Introduce
\[
 g(0)=0,\qquad g(r)=b(r-1)\ (r>0),\qquad
 G(t)=\sum_{r=0}^t g(r)=B(t-1).
\]
All prefix sums in this subsection are inclusive unless an upper
condition $y<k$ is written. The proposed lower bound for $G$ is
\begin{equation}
 p(0),\ldots,p(6)=(0,-1,0,-2,-2,-2,-2),
 \qquad p(t)=7\quad(t\geq7).
 \label{eq:tail-p}
\end{equation}

On successive complete Cantor diagonals, the sequence of second
coordinates is
\[
 (0),\quad(0,1),\quad(0,1,2),\quad\ldots.
\]
Consequently, if $t=\Delta_d+e$ with $0\leq e\leq d$, then
\begin{equation}
 H(t):=\sum_{y=0}^t g(f(y))
       =\sum_{j=0}^{d-1}G(j)+G(e).
 \label{eq:tail-first-transform}
\end{equation}
Whenever the relevant inequalities $G(j)\geq p(j)$ are available, this
gives
\[
 H(t)\geq q(t),\qquad
 q(\Delta_d+e):=\sum_{j=0}^{d-1}p(j)+p(e).
\]

The negative part of this explicit auxiliary sequence is small.
The following table records its complete finite calculation through
$t=44$, or $d=8$:
\[
\begin{array}{c|r|r}
 d&\min_{0\leq e\leq d}q(\Delta_d+e)&
 \displaystyle\sum_{e=0}^d\min(0,q(\Delta_d+e))\\ \hline
 0&0&0\\
 1&-1&-1\\
 2&-2&-4\\
 3&-3&-7\\
 4&-5&-20\\
 5&-7&-37\\
 6&-9&-58\\
 7&-11&-74\\
 8&-4&-23
\end{array}
\]
For $d\geq9$,
\[
 \sum_{j=0}^{d-1}p(j)=7d-58\geq5,\qquad p(e)\geq-2,
\]
so $q(t)\geq3$ for $t\geq45$. We have therefore proved, for all $t$,
\begin{equation}
 q(t)\geq-11,\qquad
 \sum_{t\geq0}\min(0,q(t))=-224.
 \label{eq:tail-q-negative}
\end{equation}

Apply the same diagonal identity a second time. For $k>0$, write
$k-1=\Delta_d+e$. Then
\begin{align}
 S(k)&:=\sum_{y=0}^{k-1}g(f(f(y)))\notag\\
     &=\sum_{j=0}^{d-1}H(j)+H(e)\notag\\
     &\geq\sum_{j=0}^{d-1}q(j)+q(e)
      \geq-224-11=-235.
 \label{eq:tail-235}
\end{align}
For $k=0$ the same lower bound is immediate.

Only finite prefix hypotheses are needed for this reasoning. If
$G(t)\geq p(t)$ is known for $0\leq t\leq M$, every displayed instance
with $k\leq M+1$ follows: every argument introduced by the two diagonal
decompositions is at most the original argument. This finite observation
prevents any circular use of the desired global estimate.

\subsection{The induction over all later diagonals}

We now prove Theorem~\ref{thm:tail-exact-debt}. Suppose $s\geq36$ and
the conclusions have been proved before diagonal $s$.
In particular $G(t)\geq p(t)$ is known for $0\leq t\leq s$:
its largest required earlier index is $s-1<\Delta_s$.
Thus \eqref{eq:tail-235} applies throughout the current diagonal.

Its first $s-2$ candidates, at $y=0,\ldots,s-3$, are AX words.
After the first $k$ complete AX words their relative balance is
\begin{equation}
 A(k)=10k+\sum_{y=0}^{k-1}b(f(f(y))-1)
     =10k+S(k)\geq10k-235,
 \label{eq:tail-AX-complete}
\end{equation}
where the summand is zero when $f(f(y))=0$.
Within the next AX word, its initial command adds ten and its suffix
has minimum at least $-f(f(k))\geq-k$ by
\eqref{eq:tail-suffix-loss}. For $k\geq26$, the balance after that
initial command and at every later prefix of its suffix is at least
\[
 10k-235+10-k=9k+10-235\geq0.
\]
The balance before the initial command is also nonnegative because
$10k-235\geq0$.

For the remaining finite set $0\leq k\leq25$, direct evaluation using
\eqref{eq:tail-recursions} verifies both
\begin{equation}
 A(k)\geq0,\qquad
 A(k)+10+
 \begin{cases}
   m(f(f(k))-1)&f(f(k))>0,\\
   0&f(f(k))=0
 \end{cases}
 \geq0.
 \label{eq:tail-short-AX}
\end{equation}
These are the 26 short AX checks in the independent finite-base
checker. They involve only earlier candidate values. Hence every
command prefix in the entire AX portion is nonnegative relative to
the beginning of its diagonal.

The last three candidates on diagonal $s$ are, in order,
MP at $y=s-2$, DROP at $y=s-1$, and NIP at $y=s$.
Their conservative losses, including every internal prefix, are at most
\[
 s-2,\qquad s-2,\qquad s+1,
\]
respectively. For DROP, its own $+1$ offsets one unit of the possible
loss of its suffix of length at most $s-1$. For NIP, its initial
$-1$ is added to a possible suffix loss of $s$. These bounds include
the complete words as well as their prefixes, so they may be added.
The total loss of the final three words is at most $3s-3$.

After the AX portion, \eqref{eq:tail-AX-complete} gives balance at least
$10(s-2)-235$. Every subsequent prefix is therefore at least
\[
 10(s-2)-235-(3s-3)=7s-252\geq0
\]
relative to the start of the diagonal. Thus no prefix of this diagonal
falls below its initial global balance. That initial balance is at
least seven by the induction hypothesis. This proves both the global
prefix bound and $B(n)\geq7$ on the entire new diagonal.
Lemma~\ref{lem:tail-finite-base} starts the induction.

Finally, $\mathcal W(0)$ is one NIP and has balance $-1$.
Removing it raises every later prefix balance by one. The bound becomes
$-2$, and the first NIP in candidate 5 attains it, as observed above.
This completes the proof of Theorem~\ref{thm:tail-exact-debt}.

\subsection{Soundness above an arbitrary retained tail}

Let $\mathcal D$ denote the inherited calibrated numerical calculus.
It contains truth code $0$, is closed under numerical modus ponens and
generalization, contains the eight logical rows emitted by AX, and
contains its gated theory bundle whenever the gate succeeds.
A failed gate returns truth. No assumption of consistency is made.

For a finite list $J$ of numbers and an arbitrary natural tail $S$,
define
\[
 \operatorname{push}([],S)=S,\qquad
 \operatorname{push}(x::J,S)=P(x,\operatorname{push}(J,S)).
\]
This is a representation invariant, not a claim that numerical stack
codes have a unique finite length.

\begin{theorem}[Two initial truth entries suffice]
\label{thm:tail-source-sound}
At a genuine current-source root, suppose $N=1,L=0$ and the selector and
scratch registers are zero. Suppose $T$ and $H$ are $\mathcal D$-theorems,
$D$ is arbitrary, and
\[
 W=P(0,P(0,S))
\]
for an arbitrary natural number $S$.
Every nonhalting command preserves theoremhood of the current formula
and cache, and the execution never reads $S$ as a stack.
If it halts, $\mathcal D$ derives $1$.
\end{theorem}

\begin{proof}
Track the stack as $\operatorname{push}(J,S)$, initially $J=[0,0]$,
with every member of $J$ a theorem. DROP pushes the theorem-valued
current formula and changes the top to the theorem-valued cache.
AX pushes ten theorem-valued entries and returns a theorem-valued
bundle or truth, with cache zero. An unsuccessful MP returns truth;
a successful MP obtains its new top by numerical modus ponens and
pushes an additional truth entry.

Before any possible pop, the length of $J$ equals two plus the
conservative balance so far plus the number of successful MP commands.
Theorem~\ref{thm:tail-exact-debt} bounds the conservative balance after
the prospective NIP pop below by $-2$. Therefore its value before that
pop is at least $-1$, and $J$ is nonempty. A nonhalting NIP consequently
reads a theorem from $J$, never from $S$. The identity
$\operatorname{unpair}(P(x,z))=(x,z)$ justifies this step even at the
alias $P(0,0)=0$: a conceptual leading zero is tracked by the invariant
even when its numerical insertion leaves the stack code unchanged.

If NIP halts before its pop, its incoming cache equals $1$.
That cache is a theorem by the same induction, so a native refutation
exists. There is no appeal to the theoremhood of the arbitrary tail.
\end{proof}

\subsection{Completeness and the actual startup}

For completeness, the current command compiler supplies a strong
finite block for each numerical derivation of a code $X$. From every
tuple $(t,w,0,0)$ of top, stack, counter and cache, that block either
halts earlier or returns $(X,w,0,0)$. The retained tail $w$ is arbitrary
in this contract; it need not be a stack of theorems.

For a derivation of $1$, append $\mathrm{DROP};\mathrm{NIP};
\mathrm{NIP}$ to its strong block. If the block has not already halted,
these three commands push the newly derived $1$, read it into the cache,
and halt. Prepend
\[
 \mathrm{START}_Y=\mathrm{AX}(c=0;x=1,y=Y).
\]
This command clears the counter and cache regardless of their incoming
values, so the strong block applies above whichever stack it leaves.
If $R$ is the continuation encoding of the fixed suffix, the initial
candidate code for this entire successful word is exactly
\begin{equation}
 n_Y=P(3,P(Y,P(1,R)))\geq Y.
 \label{eq:tail-cofinal}
\end{equation}
AX fields are ordered $[c+3,y,x]$. Nonempty continuations have an outer
successor, whereas the first command of a freshly selected candidate
does not; these conventions are the reason that
\eqref{eq:tail-cofinal} also covers the zero NIP tag.

Every earlier candidate drains after finitely many commands unless it
halts, by the strict descent of the suffix index. Arithmetic command
macros terminate and preserve the next-candidate counter. Thus
\eqref{eq:tail-cofinal} ensures that any refutation eventually causes
halting, from every finite cutoff and every retained working tuple.

\begin{corollary}[Calibration of the new seed]
\label{cor:tail-seed-calibration}
The current source started at
$(1,0,5,0,0,0,0,0,0)$ halts if and only if $\mathcal D$ derives $1$.
Through the inherited numerical-calculus calibration, its nonhalting
is equivalent to $\operatorname{Con}(\mathrm{ZFC})$, by ordinary
reasoning formalizable in ZFC.
\end{corollary}

\begin{proof}
Equation~\eqref{eq:tail-five} and Theorem~\ref{thm:tail-source-sound}
give soundness with $S=1$. The arbitrary-tail strong compiler and
\eqref{eq:tail-cofinal} give completeness. The inherited calibration
identifies numerical refutability with a finite ZFC refutation.
\end{proof}

This section establishes the source-level initial-condition theorem.
The complete binary runtime theorem, the exact finite blank startup,
and their binding to the released transition table remain separate
parts of the machine proof. The finite calculations here prove the
stated induction base and constants; the diagonal induction supplies
the assertion for all future candidates.
